\begin{claim}{D.2.35}{Proved}{Finite Postnikov data cannot give a universal cone-preservation invariant}
Put
\[
 B=K(\FF_2,2),\qquad A=Q\Sigma B,\qquad Y=\Sigma^2B.
\]
Every finite Postnikov section of $Y$ is ordinarily $A$-acyclic:
\[
                    P_A(P_NY)\simeq *\qquad (N<\infty).
                    \tag{D.2.35a}
\]
Consequently there is no class $\mathcal C$ of pointed CW spaces
with all of the following properties, for any fixed finite $N$:
\begin{enumerate}
\item Membership in $\mathcal C$ depends only on the $N$-type;
      $Y\in\mathcal C$ and $*\notin\mathcal C$.
\item If $C\in\mathcal C$ and
      $h:\bigvee_i\Sigma^{k_i}A\to C$, with a small indexing
      set and $k_i\geq0$, then $\cofib(h)\in\mathcal C$.
\item The class is closed under transfinite CW colimits.
\end{enumerate}
The same assertion holds for spaces with compatible marked
classes, if the marks are transported by $N$-equivalences and
by the specified successor and limit constructions.

In particular, $N\geq11$ already preserves the full mod-two
Hopf algebra and Steenrod data of $\Omega Y$ through degree ten.
A criterion involving only those data, including the primitive
operation
\[
 \operatorname{Sq}_*^1\operatorname{Sq}_*^2:
       \operatorname{Prim}H_{10}(\Omega C;\FF_2)
                    \longrightarrow H_7(\Omega C;\FF_2),
                    \tag{D.2.35b}
\]
cannot by itself supply the three properties above. This is
not a claim that $P_NY$ occurs in the particular nullification
sequence beginning at $Y$.
\end{claim}
\begin{proof}
The space $A$ is two-connected and $H_3(A;\ZZ)=\FF_2$.
Bousfield's symmetric-product comparison says that $A$ kills
$SP^\infty A$ \cite[Corollary~3.2(ii), p.~39]{Bousfield1996ULP}.
The bottom factor $K(\FF_2,3)$ is a homotopy retract of this
generalized Eilenberg--MacLane space. Thus
\[
                         P_AK(\FF_2,3)\simeq *.
                         \tag{D.2.35c}
\]
Here the bottom homology calculation follows from the
stabilization unit and the Hurewicz theorem; the factorization
of the infinite symmetric product is the usual Dold--Thom
generalized Eilenberg--MacLane decomposition.

The integral reduced homology of $Y$ is two-primary, and $Y$
is simply connected and of finite type. The Serre-class theorem
therefore makes every $\pi_iY$ a finite two-primary group
\cite[Chapter~III, Theorem~1 and Corollaries~1--2,
pp.~272--274]{I_SerreClasses1953}. In particular, $P_NY$ is a
three-connected Postnikov space with two-primary homotopy
groups. It is a generalized polyGEM by
\cite[Proposition~8.5, p.~1340]{EH5Bousfield1997}.
Apply \cite[Lemma~8.13, p.~1341]{EH5Bousfield1997} with
$G=\FF_2$, $p=2$, and $n=3$. Its extra condition applies
only when $G/pG=0$, so it imposes no condition here. The lemma
gives
\[
                    P_{K(\FF_2,3)}(P_NY)\simeq *.
                    \tag{D.2.35d}
\]
Transitivity of ordinary acyclicity, together with
\textup{(D.2.35c)}, proves \textup{(D.2.35a)}. Equivalently,
every $A$-null target is $K(\FF_2,3)$-null, and hence receives
only a contractible pointed mapping space from $P_NY$.

For completeness, the closure assertion uses actual cone
attachments, rather than an inverse-limit assertion about
nullification. Ordinary $A$-nullification can be constructed
by iteratively attaching the mapping cones of all maps from
$A$ and its nonnegative suspensions, continuing to a
sufficiently long transfinite CW colimit
\cite[Theorem~1.1 and the following construction, p.~34]
     {Bousfield1996ULP}.
At each stage the maps form a set. One may use a regular
cardinal bounding a CW model of $A$ so that maps from every
$\Sigma^kA$ into the final colimit factor through an earlier
stage up to homotopy. The final space is $A$-null, and every
cone attachment is an $A$-equivalence, which verifies the
nullification universal property.

If $\mathcal C$ had the three stated properties, the equality
of $N$-types would put $P_NY$ in $\mathcal C$. Apply this cone
construction to $P_NY$. Successor closure and transfinite
closure would put $P_A(P_NY)$ in $\mathcal C$, contradicting
\textup{(D.2.35a)} and $*\notin\mathcal C$. Compatible marked
classes give the same contradiction when the property requires
their nonvanishing.

Finally, $Y\to P_NY$ is $(N+1)$-connected, and its loop map
is $N$-connected. It consequently induces homology
isomorphisms in every degree strictly below $N$. For $N\geq11$
these identify all products, coproducts, primitives, and
Steenrod operations occurring in \textup{(D.2.35b)}, as well
as the old degree-three and degree-four loop classes and their
degree-seven commutator. Thus a preservation condition based
only on this fixed data has exactly the same initial value
on $Y$ and $P_NY$. Additional global information, or a proof
restricted to the actual intermediate targets reached from
$Y$, is necessary.
\end{proof}
